\documentclass[11pt]{article}

\usepackage[T1]{fontenc}
\usepackage[utf8]{inputenc}
\usepackage{lmodern}
\usepackage{amsmath,amssymb,amsthm,mathtools}
\usepackage[a4paper,margin=28mm]{geometry}
\usepackage{microtype}
\usepackage{xurl}
\usepackage[hidelinks]{hyperref}
\usepackage{enumitem}
\usepackage{booktabs}
\usepackage{array}

\newtheorem{theorem}{Theorem}[section]
\newtheorem{lemma}[theorem]{Lemma}
\newtheorem{proposition}[theorem]{Proposition}
\newtheorem{corollary}[theorem]{Corollary}
\newtheorem{definition}[theorem]{Definition}
\newtheorem{remark}[theorem]{Remark}

\newcommand{\N}{\mathbb N}
\newcommand{\WIN}{\ensuremath{\mathrm{WIN}}}
\newcommand{\LOSS}{\ensuremath{\mathrm{LOSS}}}
\newcommand{\DRAW}{\ensuremath{\mathrm{DRAW}}}
\newcommand{\odd}{\operatorname{odd}}
\newcommand{\src}{\operatorname{src}}
\newcommand{\ct}{\operatorname{ct}}
\newcommand{\asl}{\operatorname{asl}}
\newcommand{\Moves}{\operatorname{moves}}

\title{The Two-Player $3n\mathbin{\pm}1$ Game:\
A Binary Normal Form and a Well-Founded No-Draw Proof}
\author{Grisha Pochuev\\
  \small\href{mailto:n_854@mail.ru}{n\_854@mail.ru}\\
  \small\url{https://github.com/Grisha-Pochuev/3n-plus-minus-1-game}}
\date{10 August 2026}

\begin{document}
\maketitle

\begin{abstract}
Consider the impartial two-player game on positive odd integers in which a
move from $n>1$ chooses $3n-1$ or $3n+1$ and removes all powers of two; reaching
$1$ wins immediately. Arbitrary play need not terminate: for example
$5\to 7\to 5$ is a legal cycle. The question is whether optimal play nevertheless
has finite remoteness from every starting position. We conjugate the game to a
binary normal form with an expanding move $A$ and a strictly contracting move
$B$, introduce constant-tail coordinates and a three-free source map, and rank
marked proof configurations by retained source, finite proof-height tokens, and
a typed fixed-fibre normalizer. The key entry step is a factor-removal lemma
that either removes one factor of three at the same tail exponent or produces
finite outcome provenance. This prevents an unranked return from a factor-free
three/two frame to the preceding exponent-one state. A one-shot initialization
then allows the typed normalizer to be entered without treating a larger
temporary inner source as a source decrease. The resulting well-founded marked
normalization excludes all draw positions.
\end{abstract}

\tableofcontents

\section{The game and the theorem}

The current state is a positive odd integer $n$. For $n>1$ the player to move
chooses one of
\[
  3n-1,\qquad 3n+1,
\]
and divides the chosen value by the largest possible power of $2$, leaving an
odd integer. If the resulting odd integer is $1$, the player who made the move
wins immediately. Players alternate with perfect information.

A position is called \WIN\ for the player to move if at least one move leads to
\LOSS. It is called \LOSS\ if every legal move leads to \WIN. These two classes
are the inductively generated finite-remoteness classes. Any position not in
either class is called \DRAW.

\begin{theorem}[Main theorem]\label{thm:main}
Every positive odd starting position has finite remoteness. Equivalently, the
$3n\pm1$ game has no \DRAW\ positions, and optimal play reaches $1$ after
finitely many moves.
\end{theorem}

The theorem is not a statement about arbitrary play. There are legal cycles,
such as
\[
  5\longrightarrow 7\longrightarrow 5,
\]
under suitable choices by the players.

\section{Binary conjugation}

Write an odd state as $n=2m+1$ and define
\[
  F(m)=\left\lceil\frac{3m}{2}\right\rceil .
\]
For a nonnegative integer $x$, let $R(x)$ be the integer obtained by deleting
the maximal alternating suffix of the binary expansion of $x$. Thus
$R(1001_2)=10_2$ and $R(10101_2)=0$.

\begin{proposition}[First normal form]\label{prop:first-normal}
For $m>0$, the two children in $m$-coordinates are exactly
\[
  F(m),\qquad F(R(m)).
\]
\end{proposition}

\begin{proof}
One has
\[
  3(2m+1)-1=2(3m+1),\qquad
  3(2m+1)+1=2(3m+2).
\]
Exactly one of $3m+1$ and $3m+2$ is odd. In the other expression, each
additional factor of two removed corresponds to crossing one additional
alternating bit at the end of the binary expansion of $m$. The process stops
at the first repeated adjacent bit, which is precisely the deletion defining
$R(m)$. Returning to $m$-coordinates gives the displayed pair.
\end{proof}

Since both children lie in the image of $F$, we represent the state $F(q)$ by
$q$. The conjugated game has terminal state $0$ and moves
\[
  A(q)=\left\lceil\frac{3q}{2}\right\rceil,
  \qquad
  B(q)=R(A(q)).
\]

\begin{proposition}[Strict contraction]\label{prop:B-contracts}
For every $q>0$,
\[
  B(q)<q.
\]
\end{proposition}

\begin{proof}
Deleting a nonempty binary suffix gives
\[
  B(q)\le \left\lfloor\frac{A(q)}2\right\rfloor
       \le \left\lfloor\frac{3q+1}{4}\right\rfloor<q
\]
for $q\ge2$, while $B(1)=0$.
\end{proof}

Thus all unbounded difficulty comes from the expanding branch $A$ and from the
unbounded alternating suffix removed by $B$.

\section{Constant-tail coordinates and coefficient sources}

For a positive odd integer $a$, an exponent $r\ge1$, and a phase
$e\in\{0,1\}$, put
\[
  Q_r^e(a)=a2^r-e.
\]
This is a binary word whose final constant block has length $r$ and consists
of the bit $e$.

\begin{lemma}[Long-tail recurrence]\label{lem:long-tail}
For every positive odd $a$, phase $e$, and $r\ge3$,
\[
  A(Q_r^e(a))=Q_{r-1}^e(3a),\qquad
  B(Q_r^e(a))=Q_{r-2}^e(3a).
\]
Moreover
\[
  B(Q_1^e(a))=B(Q_2^e(a)).
\]
\end{lemma}

\begin{proof}
For $r\ge3$, the expanding child still ends in at least two equal bits; hence
its maximal alternating suffix consists only of the final bit. This gives both
long-tail identities. At the boundary, the binary words of $3a-e$ and $6a-e$
differ by appending one bit that extends the same alternating suffix, so the
same prefix remains after deletion.
\end{proof}

Define the embedded three-free coefficient
\[
  J(s)=2A(s)+1.
\]
Every positive odd coefficient $a$ has a unique factorization
\[
  a=3^kJ(s),\qquad k\ge0.
\]
Indeed, remove all factors of $3$ and call the remaining odd integer $c$.
Then $c\not\equiv0\pmod3$ and $(c-1)/2$ is $0$ or $2$ modulo $3$.
The image of
\[
 A(2t)=3t,\qquad A(2t+1)=3t+2
\]
is exactly the set of nonnegative integers not congruent to $1$ modulo $3$,
and $A$ is injective. Hence there is a unique $s$ with
$A(s)=(c-1)/2$, equivalently $c=J(s)$. Unique factorization supplies the
unique exponent $k$. The integer $s$ is called the coefficient source of
$a$ and is denoted $\src(a)$.

\begin{lemma}[Source boundary]\label{lem:source-boundary}
For every $s>0$ there is a phase
\[
  \alpha(s)=1-\left(\left\lfloor\frac{s}{2}\right\rfloor\bmod2\right)
\]
such that
\[
  \odd(3J(s)+1-2\alpha(s))=J(A(s))
\]
with $2$-adic valuation one, while the opposite phase selects $J(B(s))$ with
valuation at least two.
\end{lemma}

\begin{proof}
Put $u=A(s)$. Then $J(s)=2u+1$ and
\[
  3J(s)+1-2e=2(3u+2-e).
\]
The bracket is odd exactly when $e=1-(u\bmod2)$. For $s=2t$ or $2t+1$ one
has $u\bmod2=t\bmod2$, so this phase is exactly $e=\alpha(s)$. If $u$ is
even, then $e=1$ and $3u+1=J(u)$; if $u$ is odd, then $e=0$ and
$3u+2=J(u)$. Thus the valuation-one child is $J(A(s))$. The opposite sign is
the other original $3n\pm1$ child of the odd state $J(s)$; under the
conjugacy of Proposition~\ref{prop:first-normal} that child is $J(B(s))$.
Since it is not the valuation-one branch, its valuation is at least two.
\end{proof}

Multiplying a constant-tail coefficient by $3$ preserves its coefficient
source. Passing through a factor-free signed boundary with coefficient $J(s)$
therefore exposes exactly one ordinary $A/B$ move on the source.

For a positive state $q$, write $\kappa(q)$ for the odd coefficient in its
canonical constant-tail coordinates.

\begin{lemma}[Coefficient pullback]\label{lem:coeff-pullback}
For every positive odd $w$,
\[
  \kappa(3w)=J(R(w)),\qquad
  \kappa(3w-1)=J(R(2w-1)).
\]
\end{lemma}

\begin{proof}
Let $m\ge1$ be the length of the maximal alternating suffix of the odd word
$w$, and put $z=R(w)$. If $z>0$, maximality says that the leading bit of the
suffix equals the last bit of $z$. Because the suffix ends in $1$, its value is
\[
 t_m=\begin{cases}
 (2^m-1)/3,&m\text{ even},\\
 (2^{m+1}-1)/3,&m\text{ odd}.
 \end{cases}
\]
Thus $w=2^mz+t_m$. If $m$ is even, $z$ is even and
\[
  3w+1=2^m(3z+1)=2^mJ(z).
\]
If $m$ is odd, $z$ is odd and
\[
  3w+1=2^m(3z+2)=2^mJ(z).
\]
If $z=0$, the whole word is alternating; then $m$ is odd and
$3w+1=2^{m+1}$, whose odd part is $J(0)=1$. Hence in all cases
$\kappa(3w)=\odd(3w+1)=J(R(w))$.

For the second identity set $v=2w-1$, which is odd. Since
$3v+1=2(3w-1)$, the first identity applied to $v$ gives
\[
 \kappa(3w-1)=\odd(3w-1)=\odd(3v+1)=J(R(v))=J(R(2w-1)).
\]
\end{proof}


\section{Finite outcomes, witnesses, and ranks}

The terminal state $0$ is \LOSS\ with height $h(0)=0$. A finite \WIN\ state has
height
\[
  h(x)=1+\min\{h(y):y\text{ is a \LOSS\ child of }x\},
\]
and a finite \LOSS\ state has height
\[
  h(x)=1+\max\{h(y):y\text{ is a child of }x\}.
\]
A state whose outcome has been proved but which is being carried only to
justify a later incidence is called a \emph{routing witness}. A
\emph{retained proof token} is a finite state for which the proof has also
established the comparison needed by the rank: it either seeds an empty inner
token set at a one-shot seed transition, or replaces a previously retained
token by certified descendants of strictly smaller height. This distinction is
important: a newly discovered finite state is never silently appended to a
ranked multiset.

\begin{lemma}[Multiset token rank]\label{lem:token-rank}
For a finite multiset $M$ of proof heights define
\[
  \Phi(M)=\mathop{\#}_{h\in M}\omega^h,
\]
the natural ordinal sum of the monomials $\omega^h$. Replacing one token of
height $H$ by any finite multiset of certified descendants of heights strictly
below $H$ strictly decreases $\Phi$.
\end{lemma}

\begin{proof}
In Cantor normal form the coefficient of $\omega^H$ decreases by one, while
all inserted terms have lower exponent. No finite collection of lower terms
can compensate for the loss at exponent $H$.
\end{proof}

Every positive state $q$ has unique canonical constant-tail coordinates
$q=Q_r^e(a)$ with $a$ odd and $r\ge1$. If
$a=3^kJ(s)$, define the \emph{state source} $\rho(q)=s$.

\begin{lemma}[Universal state-source bound]\label{lem:source-bound}
For every $q>0$,
\[
  \rho(q)\le \frac{q-1}{6}.
\]
\end{lemma}

\begin{proof}
Write $q=Q_r^e(3^kJ(s))$. Since $r\ge1$ and $e\le1$,
$q\ge2J(s)-1$. Moreover
$J(s)=2\lceil3s/2\rceil+1\ge3s+1$. Hence
